\SourceLeaf{44}
\fancyhead[L]{\small EINE BESONDERE METHODE ZUR BESTIMMUNG DER POLHÖHE}
\renewcommand{\Article}[1]{\par\medskip\Needspace{5\baselineskip}\pdfbookmark[1]{Art. #1}{poli1808-art-#1}\begin{center}\textbf{#1.}\end{center}\nobreak}
\pdfbookmark[0]{Eine besondere Methode zur Bestimmung der Polhöhe}{poli1808-body}
\begin{center}\large EINE BESONDERE METHODE\par\medskip
ZUR BESTIMMUNG DER POLHÖHE.\par\medskip\rule{.09\linewidth}{.3pt}\end{center}
\Article{1}
Bekanntlich beruhen die gebräuchlichsten Methoden, mit denen Beobachter auf Reisen geografische Ortsbestimmungen vornehmen, auf Sonnenbeobachtungen. Korrespondierende Höhen dienen zur Zeitbestimmung, Höhen am oder nahe dem Meridian zur Bestimmung der Polhöhe. Diese Art der Beobachtung hat mehrere Vorzüge, besonders für diejenigen, die zur Höhenmessung den Hadley-Sextanten verwenden. Mit diesem Instrument lassen sich die Höhen der Sonne viel leichter genau beobachten als die der Sterne. Auch sind die Teilstriche am Gradbogen bei Nacht schlechter zu erkennen als am Tag. Überdies ist die Rechnung mit korrespondierenden und zirkummeridianen Höhen so einfach, dass auch jemand, der in der Mathematik wenig geübt ist, sie leicht erlernen kann. --- Lassen sich wegen Wolken oder anderer Hindernisse keine korrespondierenden Höhen gewinnen, so verbindet man gewöhnlich Beobachtungen nahe der Kulmination mit anderen, die weiter von ihr entfernt liegen. Die Abweichung der Uhr von der wahren Zeit und die Polhöhe leitet man daraus meist mit indirekten Methoden ab, etwa mit der nach \textsc{Cornelis Douwes} benannten.

Zweifellos sind Sonnenbeobachtungen den nächtlichen vorzuziehen, wann immer sie möglich sind. Wir meinen jedoch, dass auch die nächtlichen Beobachtungen Beachtung verdienen. Oft klart der Himmel bei Einbruch der Nacht auf, nachdem er tagsüber bedeckt war, oder lässt wenigstens gelegentlich
\SourceLeaf{45}die Sterne erkennen. Oft könnte auch ein Reisender, der tagsüber keine Zeit für Beobachtungen hatte, weil er seine Reise fortsetzen wollte, mit einer geografischen Ortsbestimmung an seinem Übernachtungsort einen nützlichen Beitrag leisten, sofern er sich genügend Geschick in der Beobachtung der Sterne erworben hat. Hinzu kommen die Unsicherheit der Sonnenbeobachtungen in den Wintermonaten, wenn die Sonne in unseren Gegenden kaum über den Dunst am Horizont steigt, die häufigeren Störungen durch aufdringliche Zuschauer und die Erschütterungen des künstlichen Horizonts, besonders eines Quecksilberhorizonts, durch vorbeifahrende Wagen. In der Ruhe der Nacht ist man davor besser geschützt. Man muss also zugeben, dass Sternbeobachtungen es verdienen, sorgfältiger gepflegt und praktisch genutzt zu werden.

Welchen Nutzen diese Art der Beobachtung \textit{auf Seereisen} haben kann, erkennt man leicht, wenn man bedenkt, wie wichtig es ist, die Position des Schiffes nach Möglichkeit täglich festzustellen. Davon hängen nicht selten das Wohlergehen und das Leben von Menschen ab. Gewiss haben nächtliche Beobachtungen auch auf See ihre besonderen Schwierigkeiten, denn der Meereshorizont, von dem aus die Höhen zu messen sind, ist nachts schlechter zu erkennen. Doch lässt sich diese Schwierigkeit beheben, indem man den Sextanten mit einem Fernrohr größerer Öffnung versieht.
\Article{2}
Die im vorigen Artikel angeführten Gründe für Sternbeobachtungen gewinnen noch mehr Gewicht, wenn man bedenkt, dass sich auf diese Weise in einem sehr kurzen Zeitraum, etwa durch Beobachtungen innerhalb der ersten wenigen Minuten, sowohl die Polhöhe als auch die Zeit mit erstaunlicher Genauigkeit bestimmen lassen. Dazu braucht man nur die Höhen zweier beliebiger, voneinander verschiedener Sterne zu beobachten, deren Positionen als bekannt vorausgesetzt werden, und die zugehörigen Uhrzeiten zu notieren. Der tägliche Gang der Uhr muss nur so genau bekannt sein, dass sich das Zeitintervall hinreichend genau in Sternzeit umrechnen lässt. Wenn wir uns nicht irren, zählt diese Aufgabe gewiss zu den nützlichsten der nautischen Astronomie. Daher ist es recht erstaunlich, dass diejenigen, die über diesen Gegenstand geschrieben haben, sie offenbar völlig übergangen haben. Der angesehene \textsc{Kraft} zeigte in den \textit{Actis novis Acad. Petropol.}, Bd.~XIII, wie man aus den Höhen zweier Sterne die Polhöhe bestimmt. Aus einem mir unbekannten Grund schränkte er die allgemeine Aufgabe durch eine besondere Bedingung ein. Er setzte nämlich voraus, dass beide Beobachtungen im selben Augenblick
\SourceLeaf{46}angestellt wurden. Dazu braucht man also zwei Beobachter, zwei Instrumente und genaue Gleichzeitigkeit der Beobachtungen, die sich gewiss nur schwer erreichen lässt. Dabei ist diese Bedingung völlig unnütz und überflüssig. Ein gut ausgerüstetes Schiff hat gewiss eine Uhr, die ein Zeitintervall von wenigen Minuten genau messen kann. Auch kann die Bewegung des Schiffs in so kurzer Zeit keine merkliche Änderung seiner geografischen Position bewirken. Schließlich wird die Rechnung zur Lösung der Aufgabe durch diese Bedingung um nichts einfacher, wie wir gleich zeigen werden.
\Article{3}
Leicht sieht man, dass sich die Aufgabe in der vom angesehenen \textsc{Kraft} behandelten Form auch so stellen lässt. \glqq{}Gegeben sind die Lagen zweier Punkte auf einer Kugelfläche in Bezug auf einen Großkreis, hier den Äquator. Gesucht ist die Lage eines dritten Punkts, dessen Abstände von diesen beiden Punkten gegeben sind.\grqq{} Die beiden ersten Punkte sind durch die Positionen der beiden beobachteten Sterne auf der Himmelskugel gegeben, etwa durch ihre Rektaszensionen und Deklinationen. Der dritte Punkt entspricht dem Zenit des Beobachtungsorts. Seine Deklination bestimmt die Polhöhe, seine Rektaszension den kulminierenden Punkt des Äquators und damit die Sternzeit. Die unmittelbar naheliegende Lösung verwendet drei sphärische Dreiecke. Im ersten, das aus den beiden gegebenen Punkten und dem Pol des Äquators gebildet wird, sind zwei Seiten und der eingeschlossene Winkel gegeben. Daraus ergeben sich die dritte Seite und einer der beiden übrigen Winkel, etwa der am ersten Punkt. Im zweiten Dreieck, das aus den beiden gegebenen Punkten und dem gesuchten dritten Punkt gebildet wird, sind bereits alle Seiten gegeben. Daraus ergibt sich einer der Winkel, und zwar wieder der am ersten Punkt, wenn wir im ersten Dreieck den Winkel an diesem Punkt gewählt haben. Die Summe oder Differenz dieser beiden Winkel bildet den bekannten Winkel im dritten Dreieck aus dem ersten und dem dritten Punkt und dem Pol des Äquators. In diesem Dreieck sind außerdem die beiden anliegenden Seiten bekannt. Daraus ergeben sich schließlich der Winkel am Pol und der Abstand des dritten Punkts vom Pol, also das Komplement der Polhöhe. Schon die Astronomen des 16. Jahrhunderts verwendeten diese Methode. Siehe etwa \textit{Tychonis Astronomiae instauratae progymnasmata}, S.~211 ff., wo die Positionen sehr vieler Sterne, darunter auch die des neuen Sterns in der Kassiopeia, aus ihren Abständen von je zwei Sternen bekannter Position bestimmt werden.

\SourceLeaf{47}\SourceGathering{VI.~6}Zugleich ist daraus ersichtlich, dass die Aufgabe streng genommen stets zwei Lösungen zulässt. Beschreibt man nämlich auf der Kugelfläche zwei Kreise um die gegebenen Punkte als Mittelpunkte, mit Radien gleich den jeweiligen gegebenen Abständen, so schneiden sich die Kreise im Allgemeinen in zwei Punkten. Jeder der beiden Schnittpunkte kommt als gesuchter dritter Punkt infrage. In der Praxis kann daraus jedoch keine Mehrdeutigkeit entstehen. Man denke sich den Großkreis durch den ersten und zweiten Punkt, der die Kugelfläche in zwei Halbkugeln teilt. Es kann kein Zweifel bestehen, ob der dritte Punkt, der Zenit selbst, und der Nordpol in derselben oder in entgegengesetzten Halbkugeln liegen müssen. Der erste Fall muss eintreten, wenn die Rektaszension des Sterns, der bei der Beobachtung links vom anderen stand, die Rektaszension des rechts von ihm stehenden Sterns um einen Betrag übersteigt, der, falls nötig um $360$ Grad vermehrt, zwischen $0$ und $180^0$ liegt. Der zweite Fall gilt dagegen, wenn dieser Überschuss zwischen $180^0$ und $360^0$ liegt.

Gerade dieser von vornherein in der Aufgabe liegenden Mehrdeutigkeit ist die umständliche Rechnung zuzuschreiben, die die direkte Lösung erschwert. Im Folgenden werden wir jedoch zeigen, wie man sie beträchtlich abkürzt. Lässt man auch eine indirekte Lösung zu, so kann man eine recht rasche angeben. Ihre Erläuterung behalten wir uns für eine andere Gelegenheit vor.

Bisher haben wir den Sonderfall betrachtet, in dem beide Beobachtungen gleichzeitig sind. Die allgemeine Aufgabe lässt sich auf folgende Weise auf diesen Fall zurückführen. Anstelle des an zweiter Stelle beobachteten Sterns denke man sich einen fiktiven Stern mit derselben Deklination und einer Rektaszension, die um die zwischen den beiden Beobachtungen verstrichene Sternzeit vermindert ist. Dieser fiktive Stern hätte zur Zeit der ersten Beobachtung offenkundig dieselbe Höhe erreicht, die der tatsächlich beobachtete Stern zur Zeit der zweiten Beobachtung hatte. Ersetzt man diesen durch den fiktiven Stern, so kann die Rechnung auf gleichzeitige Beobachtungen gegründet werden.

Dies alles wurde aus naheliegenden, aber geometrischen Überlegungen gewonnen. Wir hoffen daher, dass eine direkte Lösung, die allein auf analytischen Rechenschritten beruht, vielen willkommen sein wird. Sie wird gewiss erneut bestätigen, dass sich aus geometrischen Überlegungen kaum etwas gewinnen lässt, das man bei geschickter Handhabung der Analysis nicht ebenso elegant herleiten kann.
\SourceLeaf{48}
\begingroup
\newcommand{\GaussVIPoliLabel}[3]{%
  \makebox[0pt][l]{[#2]}%
  \makebox[#1\linewidth][r]{$\displaystyle #3$}%
}
\newcommand{\GaussVIPoliAligned}[1]{%
  \[\makebox[\linewidth][l]{$\displaystyle\begin{aligned}#1\end{aligned}$}\]%
}
\newcommand{\GaussVIPoliSingle}[2]{%
  \[\makebox[\linewidth][l]{%
    \makebox[0pt][l]{[#1]}%
    \makebox[\linewidth][c]{$\displaystyle #2$}}\]%
}
\Article{4}
Es sei $\varphi$ die Polhöhe, $\alpha$ und $\alpha'$ seien die Rektaszensionen, $\delta$ und $\delta'$ die Deklinationen der beiden Sterne. Ferner seien $\gamma$ und $\gamma'$ die Rektaszensionen der kulminierenden Äquatorpunkte zur Zeit der ersten beziehungsweise der zweiten Beobachtung, also die Sternzeiten selbst, in Grad umgerechnet. Schließlich sei $h$ die beobachtete Höhe des ersten Sterns und $h'$ die Höhe des zweiten. Offenbar sind $\gamma-\alpha$ und $\gamma'-\alpha'$ die zu den beiden Beobachtungen gehörenden Stundenwinkel. Wir setzen $\gamma-\alpha=\lambda$ und $\gamma'-\alpha'=\lambda-\theta$. Damit ist $\theta=\alpha'-\alpha-(\gamma'-\gamma)$ bekannt, denn $\gamma'-\gamma$ ist gleich dem Zeitintervall zwischen den beiden Beobachtungen, erst in Sternzeit und dann in Grad umgerechnet. Der Lösung der Aufgabe liegen somit diese beiden Gleichungen zugrunde:
\GaussVIPoliAligned{
\GaussVIPoliLabel{.30}{1}{\sin h}&=\sin\delta\sin\varphi+\cos\delta\cos\varphi\cos\lambda\\
\GaussVIPoliLabel{.30}{2}{\sin h'}&=\sin\delta'\sin\varphi+\cos\delta'\cos\varphi\cos(\lambda-\theta)
}
Wollten wir aus diesen Gleichungen die Unbekannte $\lambda$ oder die Unbekannte $\varphi$ eliminieren, so kämen wir zu einer zu komplizierten Gleichung. Den ersten Weg nahm der angesehene \textsc{Kraft} in der oben angeführten Abhandlung. Die so gewonnene Lösung ist jedoch nach unserem Urteil noch länger und mühsamer als die, die sich unmittelbar aus den drei Dreiecken ergibt, obwohl seine Lösung nur $\varphi$ bestimmt und die Zeit außer Betracht lässt. Es ist daher besser, eine neue Unbekannte einzuführen und damit beide Größen $\varphi,\lambda$ zu eliminieren. Bei ihrer Wahl leitet uns die folgende Beobachtung. Es gilt
\[
\begin{aligned}
&(\sin\delta\sin\varphi+\cos\delta\cos\varphi\cos\lambda)^2+(\cos\delta\sin\varphi-\sin\delta\cos\varphi\cos\lambda)^2\\
&\qquad=\sin\varphi^2+\cos\varphi^2\cos\lambda^2=1-\cos\varphi^2\sin\lambda^2
\end{aligned}
\]
und somit nach Gleichung $1$
\[
(\cos\delta\sin\varphi-\sin\delta\cos\varphi\cos\lambda)^2+\cos\varphi^2\sin\lambda^2=\cos h^2,
\]
oder
\[
\left(\frac{\cos\delta\sin\varphi-\sin\delta\cos\varphi\cos\lambda}{\cos h}\right)^2+\left(\frac{\cos\varphi\sin\lambda}{\cos h}\right)^2=1
\]

Wir können daher setzen
\GaussVIPoliAligned{
\GaussVIPoliLabel{.56}{3}{\dfrac{\cos\delta\sin\varphi-\sin\delta\cos\varphi\cos\lambda}{\cos h}}&=\cos u\\[4pt]
\GaussVIPoliLabel{.56}{4}{\dfrac{\cos\varphi\sin\lambda}{\cos h}}&=\sin u
}

\SourceLeaf{49}\SourceGathering{6\textsuperscript{*}}Kombinieren wir nun Gleichung $3$ mit $1$, so erhalten wir
\GaussVIPoliAligned{
\GaussVIPoliLabel{.34}{5}{\sin\varphi}&=\sin\delta\sin h+\cos\delta\cos h\cos u\\
\GaussVIPoliLabel{.34}{6}{\cos\varphi\cos\lambda}&=\cos\delta\sin h-\sin\delta\cos h\cos u
}

Wir bringen Gleichung $2$ in die Form
\[
\sin h'=\sin\delta'\sin\varphi+\cos\theta\cos\delta'\cos\varphi\cos\lambda+\sin\theta\cos\delta'\cos\varphi\sin\lambda
\]
und setzen für $\sin\varphi$, $\cos\varphi\cos\lambda$, $\cos\varphi\sin\lambda$ ihre Werte aus den Gleichungen $5,\ 6,\ 4$ ein. Das ergibt
\[
\begin{aligned}
&\sin h'-\sin h\sin\delta\sin\delta'-\sin h\cos\theta\cos\delta\cos\delta'-\cos u\cos h\cos\delta\sin\delta'\\
&\quad+\cos u\cos h\cos\theta\sin\delta\cos\delta'-\sin u\cos h\sin\theta\cos\delta'=0
\end{aligned}
\]

Setzen wir
\GaussVIPoliSingle{7}{\frac{\cos\delta\sin\delta'-\cos\theta\sin\delta\cos\delta'}{\sin\theta\cos\delta'}=\operatorname{cotang}v}
so geht diese Gleichung in die folgende über
\[
\begin{aligned}
&\sin h'-\sin h\sin\delta\sin\delta'-\sin h\cos\theta\cos\delta\cos\delta'\\
&\qquad-\cos h\sin\theta\cos\delta'(\cos u\operatorname{cotang}v+\sin u)=0
\end{aligned}
\]

Setzen wir also $v-u=w$, so erhalten wir
\GaussVIPoliSingle{8}{\cos w=\frac{\sin v(\sin h'-\sin h\sin\delta\sin\delta'-\sin h\cos\theta\cos\delta\cos\delta')}{\cos h\sin\theta\cos\delta'}}
woraus sich $w$ und damit $u=v-w$ bestimmen lässt.

Sobald der Winkel $u$ gefunden ist, ergibt sich $\lambda$ durch Kombination der Gleichungen $4$ und $6$, also
\GaussVIPoliSingle{9}{\operatorname{tang}\lambda=\frac{\cos h\sin u}{\cos\delta\sin h-\sin\delta\cos h\cos u}}

Aus $\lambda$ erhalten wir $\gamma=\alpha+\lambda$. Rechnen wir diesen Winkel in Zeit um, so ergibt sich der Stand der Uhr. Schließlich liefert die Kombination der Gleichungen $4$ und $5$
\GaussVIPoliSingle{10}{\operatorname{tang}\varphi=\frac{\sin\lambda\,(\sin\delta\sin h+\cos\delta\cos h\cos u)}{\cos h\sin u}}
Zur Kontrolle der Rechnung kann man bei Bedarf neben dieser Formel auch Gleichung $4$ selbst verwenden.
\SourceLeaf{50}\Article{5}
Zur Lösung aus dem vorigen Artikel fügen wir noch einige Bemerkungen hinzu.

Aus Formel $8$ ergeben sich für den durch seinen Kosinus bestimmten Winkel $w$ zwei Werte, ein positiver und ein negativer. Daraus gehen offenbar zwei verschiedene Lösungen hervor, wie wir oben bereits bemerkt haben. Welche davon die wahre ist, lässt sich jedoch leicht auf folgende Weise erkennen. Man kann zeigen, dass der Sinus der Differenz zwischen dem Azimut des Vertikalkreises, in dem der erste Stern beobachtet wurde, und dem Azimut des Vertikalkreises, in dem der zweite Stern beobachtet wurde, oder genauer der Sinus des von links nach rechts gemessenen Überschusses des ersten Azimuts über das zweite gleich dem folgenden Ausdruck ist
\[
=\frac{\sin\theta\cos\delta'\sin w}{\cos h'\sin v}
\]
Da nun $\cos\delta',\ \cos h'$ ihrer Natur nach positive Größen sind, muss $\sin w$ offenbar dasselbe oder das entgegengesetzte Vorzeichen wie $\frac{\sin\theta}{\sin v}$ haben, je nachdem, ob der erste Vertikalkreis rechts oder links vom zweiten lag. Darüber kann keine Ungewissheit bestehen, denn Beobachtungen in nahezu zusammenfallenden oder nahezu entgegengesetzten Vertikalkreisen sind für die Praxis ganz ungeeignet und daher sorgfältig zu vermeiden, wie wir unten ausführlicher darlegen werden.

Durch Einführung zweier Hilfswinkel lässt sich die Lösung noch etwas abkürzen und für die numerische Rechnung geeigneter machen. Bestimmt man nämlich den Winkel $F$ durch die Gleichung
\GaussVIPoliSingle{11}{\operatorname{tang}F=\frac{\operatorname{tang}\delta'}{\cos\theta}}
so gehen die Gleichungen $7$ und $8$ in die folgenden über
\GaussVIPoliAligned{
\GaussVIPoliLabel{.34}{12}{\operatorname{tang}v}&=\dfrac{\cos F\operatorname{tang}\theta}{\sin(F-\delta)}\\[5pt]
\GaussVIPoliLabel{.34}{13}{\cos w}&=\dfrac{\cos v\operatorname{tang}h}{\operatorname{tang}(F-\delta)}\left(\dfrac{\sin h'\sin F}{\sin h\sin\delta'\cos(F-\delta)}-1\right)
}

Führen wir ebenso einen Hilfswinkel $G$ ein, für den gilt
\GaussVIPoliSingle{14}{\operatorname{tang}G=\frac{\operatorname{tang}h}{\cos u}}
so gehen die Gleichungen $9$ und $10$ in die folgenden über
\GaussVIPoliAligned{
\GaussVIPoliLabel{.34}{15}{\operatorname{tang}\lambda}&=\dfrac{\cos G\operatorname{tang}u}{\sin(G-\delta)}\\[5pt]
\GaussVIPoliLabel{.34}{16}{\operatorname{tang}\varphi}&=\cos\lambda\operatorname{cotang}(G-\delta)
}
\SourceLeaf{51}\Article{6}
Praktische Astronomen beklagen sich öfter darüber, dass die von Theoretikern vorgeschlagenen Methoden mitunter für den praktischen Gebrauch ungeeignet sind oder zumindest hinter den Erwartungen zurückbleiben. Um einen solchen Verdacht von unserer Methode möglichst gut abzuwenden und zugleich die Anwendung der vorangegangenen Anweisungen zu zeigen, empfiehlt es sich, ein Beispiel anhand tatsächlicher Beobachtungen ausführlicher durchzurechnen. Am $21$. August dieses Jahres beobachtete ich auf unserer Sternwarte mit einem Troughton-Sextanten mit einem Radius von zehn Zoll die unkorrigierte doppelte Höhe des Sterns $\alpha$ Aquilae. Sie betrug $91^0\ 34'\ 10''$, als die Shelton-Uhr $20^{\mathrm h}\ 40^{\mathrm m}\ 8^{\mathrm s}$ anzeigte. Der Stern $\alpha$ Andromedae erreichte dieselbe Höhe bei der Uhrzeit $20^{\mathrm h}\ 46^{\mathrm m}\ 59^{\mathrm s}$. Der erste Stern wurde im rechten Vertikalkreis beobachtet, der zweite im linken. Der tägliche Gang der Uhr stimmte genau mit der wahren Sternzeit überein. Befreit man den beobachteten Winkel vom Indexfehler, so ergibt sich der wahre Winkel $91^0\ 31'\ 36''$ und damit die scheinbare Höhe $45^0\ 45'\ 48''$, aus der sich die wahre Höhe $45^0\ 44'\ 52''6=h=h'$ ergibt. Unter gehöriger Berücksichtigung von Aberration und Nutation ergaben sich folgende Sternpositionen:
\[
\begin{array}{rcl@{\quad}rcl}
\alpha&=&295^0\ 22'\ 6''6&\delta&=&+8^0\ 22'\ 43''1\\
\alpha'&=&359\ 38\ 18.5&\delta'&=&+28\ 2\ 13.4
\end{array}
\]
Die in einen Bogen umgerechnete Zeitdifferenz $6'\ 51''$ ergibt $\gamma'-\gamma=1^0\ 42'\ 45''$, woraus $\theta=62^0\ 33'\ 26''9$ folgt. Die Rechnung selbst verläuft nun folgendermaßen (der Faktor $\frac{\sin h'}{\sin h}$ in Formel $13$ fällt hier wegen der gleichen Höhen offenbar weg).
\[
\begin{array}{@{}l@{\quad}r@{}l@{\hspace{.65em}}l@{}}
\log\operatorname{tang}\delta'\quad .\ .\ .\ .&9.7263516&&\\
\log\cos\theta\quad .\ .\ .\ .\ .&9.6635677&&\\\cline{1-3}
\log\operatorname{tang}F\quad .\ .\ .\ .&0.0627839&,&\text{daraus }F=49^0\ 7'\ 37''70\\
&&&\text{und }F-\delta=40^0\ 44'\ 54''60\\[8pt]
\log\operatorname{tang}\theta.\quad .\ .\ .\ .&0.2845877&&\\
\log\cos F\quad .\ .\ .\ .\ .&9.8158318&&\\
\text{Komp. }\log\sin(F-\delta)\quad .&0.1802598&&\\\cline{1-3}
\log\operatorname{tang}v\quad .\ .\ .\ .\ .&0.2856793&&\text{daraus }v=62^0\ 36'\ 58''79\\[8pt]
\log\sin F\quad .\ .\ .\ .\ .&9.8786157&&\\
\text{Komp. }\log\sin\delta'\quad .\ .\ .&0.3278629&&\\
\text{Komp. }\log\cos(F-\delta)\quad .&0.1205703&&\\\cline{1-3}
&0.3270489&&=\log 2.123483=\log n
\end{array}
\]
\endgroup
